\chapter{Literature Survey}\label{chap:literature_survey}

\section{Overview of Existing Techniques}

The problem of Deterministic Finite Automaton (DFA) minimization has been explored for several decades as one of the most fundamental optimization problems in automata theory. The goal of minimization is to produce an equivalent automaton with the smallest possible number of states, ensuring computational efficiency in pattern recognition, lexical analysis, and model checking tasks. Over time, research has evolved from purely sequential approaches to parallel and GPU-accelerated algorithms, reflecting the changing nature of available hardware.

Early methods such as Moore’s algorithm (1956) laid the foundation for iterative refinement of state partitions based on distinguishable transitions. However, the first major efficiency leap came with Hopcroft’s algorithm (1971), which introduced a clever selection of splitters to achieve a time complexity of O(nlogn), a result that stood as a practical benchmark for decades. Hopcroft’s refinement logic inspired numerous variants optimized for specific machine architectures and language-processing tasks.

Subsequent work focused on improving the data structures and memory access patterns used in these algorithms. The Paige–Tarjan algorithm (1987) extended the partition refinement concept and introduced methods applicable beyond DFAs, such as bisimulation and coarsest partitioning problems in process algebra and graph theory. This approach achieved near-linear performance on sparse transition systems and became a theoretical cornerstone for many later parallel formulations.

With the rise of parallel computing models in the late 20th century, researchers began exploring how DFA minimization could be adapted to exploit multi-processor and GPU architectures. While parallel computing offered the promise of faster runtime through concurrency, it also introduced challenges in synchronization, load balancing, and memory contention. Notably, Tewari, Ravikumar, and Arora (2002) examined the problem through the lens of Parallel Random Access Machine (PRAM) models, demonstrating that DFA minimization belongs to the complexity class Nick’s Class (NC), and can, in theory, be computed efficiently in parallel.

More recent developments have sought to transform these theoretical results into practical, GPU-executable algorithms. Work by Wijs (2015) on GPU-accelerated bisimilarity checking paved the way for efficient implementations of partition-based automata algorithms on many-core hardware. The base paper of this seminar continues this trajectory, evaluating multiple massively parallel DFA minimization algorithms and assessing their scalability and memory performance on contemporary GPU systems.

Together, these studies show a clear evolution—from simple sequential refinement to hybrid parallel algorithms that integrate both theoretical rigor and hardware-aware engineering.

\section{Key Studies and Methodologies}

\subsection{Hopcroft’s O(nlogn) Minimization Algorithm (1971)}

\textbf{Methodology:}
Hopcroft’s algorithm is based on iterative partition refinement of DFA states into equivalence classes. The key innovation lies in its strategy of selecting the smallest active partition block to refine at each step. By avoiding redundant re-examination of large partitions, the method ensures that each state is moved between blocks a logarithmic number of times, leading to the optimal asymptotic complexity of O(nlogn). This design drastically reduced the computational cost compared to Moore’s earlier O(n^2) approach.

\textbf{Strengths:}
\begin{itemize}
\item Excellent performance on deterministic input, even for large state sets.
\item Simple data structure design suitable for software implementations.
\item Serves as the reference baseline for evaluating later algorithms.
\end{itemize}

\textbf{Weaknesses:}
\begin{itemize}
\item Inherently sequential — difficult to parallelize because partition refinement depends on the order of processed splitters.
\item Performance deteriorates for automata with high branching factors or irregular transition distributions.
\end{itemize}

\subsection{Partition Refinement and Paige–Tarjan Framework (1987)}

\textbf{Methodology:}
Paige and Tarjan introduced a generalized partition refinement framework applicable to equivalence problems beyond DFA minimization. Their algorithm efficiently refines partitions of a relation by splitting sets based on relational dependencies, achieving near-linear performance on sparse graphs. It forms the conceptual basis for many coarsest-partition algorithms in model checking, bisimulation, and control theory.

\textbf{Strengths:}
\begin{itemize}
\item General-purpose: can be applied to both DFA minimization and bisimulation.
\item Highly efficient on large systems with sparse transitions.
\item Facilitates incremental refinement, enabling early termination when no further splits occur.
\end{itemize}

\textbf{Weaknesses:}
\begin{itemize}
\item Although asymptotically efficient, its stepwise refinement is not naturally parallelizable.
\item Synchronization between refinement steps can cause overhead in distributed environments.
\end{itemize}

\textbf{Relevance to Modern Parallel Systems:}
This framework inspired most of the parallel DFA minimization algorithms developed later. Many GPU-based approaches—including those in Wijs (2015) and the base paper—derive their internal refinement logic from the Paige–Tarjan structure, adapting it to parallel primitives such as prefix sums, hash-based partitioning, and warp-synchronous operations.

\subsection{Parallel DFA Minimization (Tewari, Ravikumar \& Arora, 2002)}

\textbf{Methodology:}
Tewari et al. presented one of the first systematic studies of DFA minimization in parallel environments. They explored multiple PRAM-based strategies to achieve polylogarithmic time complexity using a polynomial number of processors. The core technique involves representing the DFA as a state-pair graph and performing a parallel transitive closure to identify distinguishable state pairs. In theory, this allows DFA minimization in 𝑂(log^2 𝑛) parallel time using O(n^2) processors.

\textbf{Strengths:}
\begin{itemize}
\item Establishes DFA minimization as a problem within the NC complexity class.
\item Introduces transitive closure as a universal parallel primitive for equivalence detection.
\item Lays the groundwork for modern GPU-based algorithms by identifying the parallelizable components of the process.
\end{itemize}

\textbf{Weaknesses:}
\begin{itemize}
\item Requires an impractically high number of processors for large DFAs.
\item Memory consumption grows quadratically with the number of states.
\item Synchronization overhead between processors can lead to diminishing returns on physical hardware.
\end{itemize}

\textbf{Impact:}
Although the PRAM-based method is mainly of theoretical interest, it provided an essential blueprint for later hybrid algorithms that combine partition refinement with partial transitive closure steps—an approach explored in both Wijs (2015) and the 2024 base paper.

\subsection{GPU-Accelerated Bisimilarity Checking (Wijs, 2015)}

\textbf{Methodology:}
Wijs’ research focuses on adapting bisimilarity and equivalence checking algorithms for GPUs. The study utilizes CUDA to parallelize operations such as partition updates, state-pair evaluation, and block refinement. The work introduces memory-optimized data layouts, avoiding dense adjacency matrices and instead using compressed sparse representations to handle large state spaces efficiently. GPU primitives like parallel sort, reduce, and exclusive scan are used to update partitions efficiently across thousands of threads.

\textbf{Strengths:}
\begin{itemize}
\item Demonstrates the feasibility of executing large-scale state-space minimization on GPUs.
\item Provides clear evidence of speedups compared to CPU implementations, particularly for models with millions of transitions.
\item Establishes reusable GPU design principles applicable to DFA minimization (data coalescing, warp-level synchronization).
\end{itemize}

\textbf{Weaknesses:}
\begin{itemize}
\item Strong dependence on the GPU’s memory capacity; large DFAs can still exceed VRAM limits.
\item Performance varies with automata structure; irregular graphs can cause load imbalance among GPU threads.
\end{itemize}

\textbf{Relevance:}
Wijs’ methods directly inspired the experimental setup of the 2024 study, which builds upon the GPU partition-refinement model to test multiple DFA minimization algorithms under realistic workloads.

\subsection{Modern Parallel and Hybrid Approaches (Martens \& Wijs, 2024)}

\textbf{Methodology:}
The base paper evaluates three categories of massively parallel DFA minimization algorithms:

\begin{itemize}
\item Transitive-closure-based algorithms — theoretically optimal but memory intensive.
\item Partition-refinement-based methods — practical and GPU-friendly.
\item Hybrid models — combining partial closure steps with refinement to reduce iterations.
\end{itemize}

Experiments were performed on Fibonacci, Bit-Splitter, and VLTS datasets, measuring execution time, scalability, and memory efficiency. The study showed that partition-refinement techniques outperform theoretical transitive-closure algorithms in practice due to lower memory footprint and better GPU utilization.

\textbf{Strengths:}
\begin{itemize}
\item Comprehensive benchmarking across multiple datasets.
\item Offers empirical insight into how GPU hardware interacts with algorithmic structure.
\item Suggests an optimal balance between theoretical complexity and implementation feasibility.
\end{itemize}

\textbf{Weaknesses:}
\begin{itemize}
\item The approach’s performance depends heavily on dataset characteristics.
\item Hybridization introduces additional tuning parameters for closure depth and iteration control.
\end{itemize}

\section{Evolution of Parallel DFA Minimization}

The transition from sequential to parallel DFA minimization can be viewed as a three-stage evolution:

\begin{itemize}
\item Theoretical groundwork (Hopcroft, Paige–Tarjan) — defining efficient sequential procedures.
\item Parallel complexity theory (Tewari et al.) — proving DFA minimization belongs to NC, even if resource-intensive.
\item Practical GPU realization (Wijs, Martens) — designing real-world implementations that exploit massive parallelism without excessive memory cost.
\end{itemize}

The main challenge lies in bridging the gap between theoretical concurrency and hardware-level parallelism. GPUs, while powerful, impose restrictions on global synchronization and memory access patterns. Effective designs must consider factors such as warp divergence, shared memory usage, and coalesced access, as these significantly influence real-world speedups.

\section{Summary of Limitations}

From the works above and recent GPU-focused evaluations, several recurring limitations emerge:

\begin{itemize}
\item \textit{Resource blow-up for theory-optimal algorithms.} NC-style or transitive-closure based parallel approaches can give excellent theoretical depth but require enormous numbers of processors or quadratic/cubic-sized intermediate structures, making them impractical on real GPUs except for tiny inputs. Practical studies show these approaches often run out of memory or time on moderate-sized DFAs.
\item \textit{Sensitivity to input shape (worst-case automata).} Families such as Fibonacci automata and certain bit-splitter constructions force partition-refinement algorithms to perform many iterations or expensive splits; these worst-case inputs reveal gaps between average-case and pathological behavior. Empirical benchmarks used in recent evaluations specifically highlight these pathological cases and show radically different algorithmic rankings across families.
\item \textit{GPU memory and access-pattern constraints.} Many theoretically parallel operations (matrix-style reachability encodings, product-graph representations) produce dense or very large structures. GPUs have high throughput but limited global memory and specific memory-access performance characteristics; algorithms must be restructured (use sorting, compact encodings, parallel scans) to avoid memory stalls and excessive allocation.
\item \textit{Trade-off between iterations and per-iteration cost.} Sorting/signature-based methods can reduce the number of refinement iterations (by splitting blocks into many subblocks in one pass) but pay a higher per-iteration cost (sorting and signature recomputation). Conversely, leader-election / incremental split strategies have cheaper iterations but may require many more of them. Choosing between them depends strongly on the dataset distribution and alphabet size.
\item \textit{Dependence on PRAM primitives and atomics.} Several parallel algorithms assume CRCW PRAM semantics or heavy use of atomic operations (CAS). Implementing these semantics efficiently on GPUs requires careful locking/atomic strategies; sometimes atomic contention nullifies theoretical gains.
\end{itemize}